\section{Step 13 Form capital sigma one of old e}\label{sha-s13}

Apply the three word rotations to \texttt{s4}, which still holds the old e value. XOR the results into \texttt{t0}. The working vector must remain unchanged until both round temporaries are ready. This capital-sigma function has no logical-shift term; it differs from the schedule's small-sigma-one function.

\subsection{Write and inspect}
Replace only the temporary instructions between \texttt{\# BEGIN comp\_sigma1} and \texttt{\# END comp\_sigma1}. Keep the supplied labels outside that region.

\Needspace{9\baselineskip}
\MARSCodeLabel{MIPS assembly}
\begin{minted}[samepage=true]{mars-idea}
    rotr $t0, $s4, 6
    rotr $t2, $s4, 11
    xor $t0, $t0, $t2
    rotr $t2, $s4, 25
    xor $t0, $t0, $t2
\end{minted}

Use recipe \texttt{s13-1} with the cumulative program through this edit. The exact recipe names the entry, prepared input and checkpoint. Observe the checkpoint before continuing to its temporary or normal return.

\subsection{Expected observation}
For round zero of padded abc, \texttt{t0} has low word \texttt{3587272b}.

\Needspace{9\baselineskip}
\subsection{Check your understanding}
Which value supplies all three rotations?

\begin{itemize}
\item \textbf{A} The new e after its replacement.
\item \textbf{B} The old e in \texttt{s4}.
\item \textbf{C} The byte count.
\end{itemize}
